\section{Workshop Goals}\label{sec:workshopGoals}
%aims to be an inclusive forum to address four nontrivial challenges in uncertainty visualization, as identified through the 
%The proposed workshop is continuation of the {\em 2024 IEEE Uncertainty Visualization Workshop}

\vspace{-1mm} The proposed workshop aims to expand on the four nontrivial challenges identified in the {\em 2024 IEEE Uncertainty Visualization Workshop}, namely uncertainty propagation, uncertainty communication, AI, and decision-making (see Fig.~\ref{fig:teaser}), and be an inclusive forum to address the identified challenges. 
%We aim to discuss the solutions to these four identified challenges in the workshop.  

%There have been two prior VIS workshops held in years 2011 and 2015 on the topic of uncertainty led by Dr. Johnson and Dr. Potter, respectively. 
%problems in uncertainty through this workshop. Specifically, we elaborate our rationale and objectives pertinent to the four areas in uncertainty visualization: applications, techniques, software, and decision frameworks.  

%State-of-the-art literature lacks the theory regarding optimal models of uncertainty, as pointed in the keynote talk by Dr. Weiskopf in 2024 uncertainty visualization workshop.

\textbf{Uncertainty Propagation:}
Uncertainty should be considered as data in its own right since it exists in all simulation, experimental, and observational data. Currently, uncertainty is often modeled with probability distributions (Fig.\ \ref{fig:teaser}a). However, it is not clear which data models are optimal for uncertainty representation inside visualization tools. For example, is a histogram the best representation to characterize uncertainty and most easily interpretable by users? How can correlation in uncertain data be represented for the accuracy of visualization? Furthermore, linear/nonlinear transformation of models of uncertainty by complex data processing and visualization algorithms (Fig.\ \ref{fig:teaser}b) can make uncertainty propagation intractable. Thus, we propose to discuss solutions to the following critical questions related to uncertainty propagation. (1)~What are optimal models of uncertainty that are easy-to-use within visualization tools and enable capturing correlation for efficient and accurate results? (2)~What are the bottlenecks that prevent us from propagating the uncertainty? (3)~What are theoretical and application-specific approaches that enable tracking of uncertainty in complex scientific workflows? 


%Lastly, creating models of uncertainty for perception,
%cognition is even harder. For example, how good users of visualization tools are at interpreting different uncertainty models. How uncertainty propagation through
%nonlinear transformations can be tracked (through solutions other
%than Monte Carlo).
%There is currently no theory regarding optimal models of uncertainty.
%Another limitation of existing uncertainty models is that many of them assume that the data are independent due to their simplicity. However, these models can lead to inaccurate estimate of uncertainty compared to more complex models that consider data correlation, as pointed in the keynote talk in 2024 uncertainty workshop.

%One of the prime challenges identified in the 2023 Dagstuhl workshop on Decision-Making Under Uncertainty~\cite{boukhelifa_et_al:DagRep.12.8.1} includes a lack of understanding regarding intersection of uncertainty visualization and application domains. The Dagstuhl workshop identified success stories of publication of uncertainty visualization in seven application domains, including oceanology, medicine and so on (e.g., see Figure~\ref{fig:teaser}a). The three main questions that were raised based on the listed success stories are: (1) Distinguishing which applications need visualizing uncertainty and which do not, (2) What kind of uncertainty needs to be depicted depending on data and tasks, (3) How far commercial software has progressed in understanding uncertainty in respective application domains. The goal of this workshop is therefore to {\em advance understanding of how uncertainty is understood and treated in the context of application domains} by inviting submissions and through direct interaction with application scientists.
%A well-known example in uncertainty visualization community is the cone of uncertainty to convey uncertainty in hurricane tracks is often wrongly interpreted as size of storm by general users

%The contributions in effective uncertainty communication will throw light on the following three important aspects: (1) novel approaches for understanding uncertainty in 2D/3D/high-dimensional data from various application domains, (2) bottlenecks that prevent us from visualizing uncertainty in high-dimensional data, (3) ways to improve the design of existing uncertainty visualizations through expert interactions and feedback.  

\textbf{Uncertainty Communication:} As discussed by the panelists of the 2024 uncertainty workshop, the ultimate goal of communicating uncertainty is to enhance trust and decision-making among users in a consistent manner.  However, users often struggle to interpret uncertainty if it is not communicated properly. Although the use of 1D uncertainty visualization techniques (e.g., error bars, box plots) is standardized in visualization tools to convey uncertainty in data, there are no established ways to communicate uncertainty in 2D/3D/high-dimensional data, and users can feel impaired when interpreting such uncertainties. We will, therefore, invite submissions and organize sessions that will help to address the following challenges: (1) Understanding the effective ways to visually communicate uncertainty that can be interpreted well by both scientific audiences as well as non-experts (Fig.\ \ref{fig:teaser}c). (2) Understanding the pedagogical aspects of uncertainty visualization (Fig.\ \ref{fig:teaser}d) that will be key to enabling users to interpret uncertainty and make decisions effectively. In addition, our workshop will present an opportunity for researchers to publish novel methods for uncertainty communication that are not accepted or are too preliminary for the main VIS event. 

\textbf{AI and Uncertainty Visualization:} The use of AI is rapidly growing in science due to the tremendous power it brings to enhance the efficiency of data visualization and analysis. The use of AI models in science, however, is still quite limited because these models are not perfect and can reduce the credibility of predicted results. Furthermore, many existing uncertainty visualization approaches are application-specific. However, we need more general solutions for broader impact. In this context, advances in AI present a great opportunity for the generalization of uncertainty visualization. For example, federated learning can be utilized in innovative ways to visualize uncertainty across applications. We will, therefore, aim to find solutions that can help answer the following two critical questions related to AI. (1)~Can uncertainty visualization of AI models be derived to support explainable AI (Fig.\ \ref{fig:teaser}g) and increase trust in predicted results? (2)~Can AI surrogates be used to understand uncertainty propagation in complex data processing and visualization algorithms for generalizability and broader impact (Fig.\ \ref{fig:teaser}h)?  

%Can AI help to generate uncertainty models of input data? Can AI surrogates be used to understand uncertainty propagation in complex data processing and visualization algorithm? (Fig.\ \ref{fig:teaser}h) Furthermore, AI models themeselve are not perfect. Can uncertainty visualization of AI models be derived to support explainable AI (Fig.\ \ref{fig:teaser}g)? 
%and how such evaluation can be generalized for broader class of problems (e.g., scalar field uncertainty)

\textbf{Decision-Making Under Uncertainty:} Although visualization of uncertainty provides a fuller representation of data, making optimal decisions under uncertainty can be difficult. For example, users may not understand what 80\% probability of the event means and how to make decisions based on probabilities. Few recent research utilized statistical measures including, median~\cite{contourBoxplots} and mode~\cite{TA:Athawale:2016:nonparametricIsosurfaces} (Fig.~\ref{fig:teaser}e), for robust decision-making under uncertainty. However, the theoretical understanding of decision-making driven by uncertainty visualization is still in the early stages. Furthermore, it is also critical to evaluate how user perception and cognition play a role in decision-making (Fig.~\ref{fig:teaser}f).  We will, therefore, invite submissions that foster discussions on the following topics that will help us advance our understanding of optimal decision-making under uncertainty. (1)~Identification of obstacles to optimal decision-making when users are presented with uncertainty. (2)~New techniques for decision-making under uncertainty. (3)~User evaluation and generalization of existing decision frameworks for a more general class of visualization problems (e.g., scalar fields). (4)~Innovative use of AI to learn from past user decisions for optimal decision-making under uncertainty.
%Inviting submissions on methods or application use cases for resolving uncertainty therefore will be a key to developing future strategies for robust decision-making based on visualization of uncertainty.

% \vspace{-1mm}Our proposed workshop aims to be an inclusive forum for the fast dissemination of the latest results in uncertainty visualization and related concepts. There have been two VIS workshops in 2011 and 2015 on the topic of uncertainty led by Dr. Johnson and Dr. Potter, respectively. Given a significant progress in uncertainty research in the years following these workshops, we aim to discuss the solutions for persistent and new problems identified in the four areas in uncertainty visualization: applications, techniques, software, and decision frameworks.  

% %problems in uncertainty through this workshop. Specifically, we elaborate our rationale and objectives pertinent to the four areas in uncertainty visualization: applications, techniques, software, and decision frameworks.  

% \textbf{Applications:} One of the prime challenges identified in the 2023 Dagstuhl workshop on Decision-Making Under Uncertainty~\cite{boukhelifa_et_al:DagRep.12.8.1} includes a lack of understanding regarding intersection of uncertainty visualization and application domains. The Dagstuhl workshop identified success stories of publication of uncertainty visualization in seven application domains, including oceanology, medicine and so on (e.g., see Figure~\ref{fig:teaser}a). The three main questions that were raised based on the listed success stories are: (1) Distinguishing which applications need visualizing uncertainty and which do not, (2) What kind of uncertainty needs to be depicted depending on data and tasks, (3) How far commercial software has progressed in understanding uncertainty in respective application domains. The goal of this workshop is therefore to {\em advance understanding of how uncertainty is understood and treated in the context of application domains} by inviting submissions and through direct interaction with application scientists.

% \textbf{Techniques:} Techniques for uncertainty visualization of 2D/3D/high-dimensional data are not as ubiquitous 1D uncertainty visualization techniques (e.g., error bars, box plots). There have been only a few advances by the visualization community that researched how to display errors or uncertainties in 2D/3D spatial data~\cite{TA:Potter:2009:ensemble-vis, TA:Pothkow:2011:probMarchingCubes, TA:Athawale:2016:nonparametricIsosurfaces, mandatoryCriticalPoints} (Figure~\ref{fig:teaser}b), high-dimensional graph data~\cite{graphUncertainty, networkUncertainty}. We will therefore invite submissions that specifically {\em target uncertainty visualization and analysis techniques for 2D/3D/high-dimensional data.} Our workshop will be a good venue for researchers to publish  uncertainty visualization and analysis techniques that are not accepted or that are too preliminary to publish in the main VIS event. The contributions on uncertainty visualization and analysis techniques will throw light on the following three important aspects: (1) novel approaches for understanding uncertainty in 2D/3D/high-dimensional data from various application domains, (2) bottlenecks that prevent us from visualizing uncertainty in high-dimensional data, (3) ways to improve the design of existing uncertainty visualizations through expert interactions and feedback.  

% \textbf{Software:} One important aim of this workshop is to {\em progress the community understanding of what uncertainty visualization software is available (e.g., Figure~\ref{fig:teaser}c) and needed across universities, laboratories, and commercial organizations}. Existing software tools used in visualization community, e.g., ParaView~\cite{AHRENS2005717} and VisIt~\cite{VisIt2011}, provide lots of data processing filters, but lack functionalities that can convey uncertainty in results produced by these filters. Inviting submissions on uncertainty visualization software/tools will provide the visualization community two advantages: (1) assessment of the progress in uncertainty visualization software, (2) development of a roadmap for understanding what software is needed to meet application needs and promoting research collaborations.

% \textbf{Decision Frameworks:} Although visualization of uncertainty provides a fuller representation of data, interpreting uncertainty and making optimal decisions under uncertainty can be difficult. Few recent research probabilistically quantified data uncertainty and utilized nonparametric measures including, median~\cite{contourBoxplots} and most-probable~\cite{TA:Athawale:2016:nonparametricIsosurfaces} frameworks (Figure~\ref{fig:teaser}d), for robust decision-making under uncertainty. {\em The goal of this workshop is therefore to advance the understanding of decision frameworks that can be applied to resolve uncertainty in a robust manner.} Inviting submissions on methods or application use cases for resolving uncertainty therefore will be a key to developing future strategies for robust decision-making based on visualization of uncertainty.